% Proposed B.151. Owned isolated source; no central edit or private render.
\subsection*{The complete one-sided filling and its free family kernel}
\statement{B.151}{Proved}{The entire first filling has an exact free-coset kernel, while the opposite square remains coupled to the base}
\label{stmt:B151}
Retain the actual ordinary groups and both marked families of B.141.
There is a retract
\[
 \mathcal B=L_L*_KR_R\ \le\ V=R_L*_KR_R,\qquad
 L=\langle K,k\rangle,
\]
such that \(V\) is a free-coset extension of \(\mathcal B\).
All five first-family relators lie in \(\mathcal B\). If
\[
 R_i=\langle\!\langle P_0,\ldots,P_4\rangle\!\rangle_{\mathcal B},
 \quad M_i=\langle\!\langle P_0,\ldots,P_4\rangle\!\rangle_V,
 \quad \mathcal B_i=\mathcal B/R_i,
\]
then
\begin{equation}\label{eq:B151-main}
\begin{gathered}
 M_i\cap\mathcal B=R_i,\\
 V/M_i=F\left(\coprod_j\mathcal B_i/S_{ji}\right)
                    \rtimes\mathcal B_i.
\end{gathered}
\end{equation}
The four actual subgroups \(S_{ji}\) are specified below as
image subgroups, without assuming their presentations survive.
The restriction of \(M_i\) to the original free family kernel
is exactly the normal closure of the induced coset identifications.
After this complete first step, the four opposite leaves lie
in the free kernel and the opposite square retains the explicit
base component in \eqref{eq:B151-opposite-square}.
All statements are ordinary; no abstract completion of the
infinite free family kernel or final filling survival is asserted.

\begin{proof}
\emph{B.151.1. The complete regional centralizer presentation.}
Use B.145's actual subgroup and auxiliary section
\begin{equation}\label{eq:B151-L}
\begin{gathered}
 L=\langle h,k,a,d,e,V_4\mid
 [h^2,k],[k,a],[k,d],[k,e],[k,V_4]\rangle,\\
 T_h=F(h,a,d,e,V_4),\qquad
 Q_0=\langle h^2,a,d,e,V_4\rangle.
\end{gathered}
\end{equation}
The four complete family orbits in B.147 have representatives
\[
 \lambda_\alpha=\alpha_1,\quad \lambda_\beta=\beta_0,\quad
 \lambda_\delta=\delta_1,\quad \lambda_\epsilon=\epsilon_1
\]
and stabilizers
\begin{equation}\label{eq:B151-stabilizers}
\begin{aligned}
 Q_\alpha&=\langle hah^{-1},h^2V_4h^{-2},hdh^{-1},heh^{-1}\rangle,\\
 Q_\beta&=\langle a,hV_4h^{-1},hdh^{-1},heh^{-1}\rangle,\\
 Q_\delta=Q_\epsilon&=h\langle a,d,e,V_4\rangle h^{-1}.
\end{aligned}
\end{equation}
Set \(S_j=\langle k,Q_j\rangle\le L\), as actual subgroups.
The full ordinary regional group has the exact presentation
\begin{equation}\label{eq:B151-regional}
 R^0=\langle L,\lambda_j\ (j=\alpha,\beta,\delta,\epsilon)
              \mid [\lambda_j,s]=1\ (s\in S_j)\rangle.
\end{equation}
Only \(k\) and the four displayed generators of each \(Q_j\)
are needed in this finite presentation over \(L\).

Indeed B.141's entire family set of type \(j\) is \(T_h/Q_j\);
send its vertex \(tQ_j\) to \(t\lambda_jt^{-1}\).
The stabilizer relations make this independent of the representative.
Every center/family edge is a \(T_h\)-translate of the edge
\([k,\lambda_j]=1\); its full neighborhood is the \(Q_j\)-orbit
of \(k\). Commuting with both \(k\) and \(Q_j\) gives all
these edges. Conversely each displayed stabilizer and edge
relation occurs in the actual graph presentation. There are no
family/family edges. These substitutions and the root-vertex
assignments give inverse presentation maps, proving
\eqref{eq:B151-regional} with every height and retained coset.

Inside \(L\), its center-kernel coset basis also gives
\[
 S_j=F\bigl(X_{qQ_0}:q\in Q_j\bigr)\rtimes Q_j.
\]
Here the distinct center indices are
\(Q_j/(Q_j\cap Q_0)\). No relation \([k,Q_j]=1\) is being
imposed: the new family letter centralizes their generated subgroup.

\emph{B.151.2. Retaining the full opposite region.}
The actual color retraction \(R_L\to L_L\) fixes \(K\).
Ordinary amalgam normal form embeds
\(\mathcal B=L_L*_KR_R\) in \(V\), and the induced map
\(\rho:V\to\mathcal B\) is a retraction. Substitution of
\eqref{eq:B151-regional} gives
\begin{equation}\label{eq:B151-over-B}
 V=\langle\mathcal B,\lambda_j\mid[\lambda_j,S_j]=1\rangle.
\end{equation}
In particular every right-region family remains in \(\mathcal B\).

For any group \(B\) and specified subgroups \(S_j\le B\),
the presentation on the right of \eqref{eq:B151-over-B} is
\begin{equation}\label{eq:B151-coset}
 F(\Omega)\rtimes B,\qquad \Omega=\coprod_j B/S_j,
\end{equation}
with the left coset action. One map sends the free letter
\(X_{bS_j,j}\) to \(b\lambda_jb^{-1}\). Its representative
independence is exactly the centralizer relation. The inverse
fixes \(B\) and sends \(\lambda_j\) to \(X_{S_j,j}\).
They agree inversely on every generator and preserve the
specified action. Applying this to the actual \(\mathcal B\)
proves the asserted complete free family kernel.

\emph{B.151.3. The whole first normal closure.}
Write \(\mathfrak d\) for B.141's marked map denoted there by \(d\).
Thus
\begin{equation}\label{eq:B151-words}
\begin{gathered}
 \mathfrak d(u)=p\bigl(aC(a^{-1}\eta(u)a)Ca^{-1}\bigr),\\
 (V_1,V_2,V_3,V_4)=(a,d,e,f^{-1}Bf),\qquad
 \gamma=\mathfrak d(f),\quad x_i=\mathfrak d(V_i),\\
 P_0=[(\gamma k)^2,k],\qquad P_i=[k,x_i]\quad(1\le i\le4).
\end{gathered}
\end{equation}
All five are in \(\mathcal B\), since \(\mathfrak d(K)\)
is in the full right region and \(k\in L_L\).
Adding these five relations to \eqref{eq:B151-over-B}
replaces only its base by \(\mathcal B_i\) and its stabilizers
by their actual images \(S_{ji}\). The two inverse coset maps
in \eqref{eq:B151-coset} still apply. This proves
\eqref{eq:B151-main}; equivalently \(\rho\) descends to a
retraction onto \(\mathcal B_i\).

More precisely, the family-coordinate map is
\[
 F(\Omega)\longrightarrow F(R_i\backslash\Omega)
                  =F\left(\coprod_j\mathcal B_i/S_{ji}\right).
\]
Its kernel is exactly the free-group normal closure of
\[
 X_\omega X_{r\omega}^{-1}
       \qquad(r\in R_i,\ \omega\in\Omega).
\]
This follows either by freely identifying each orbit of basis
letters or directly from the two presentation maps. It evaluates
the restriction of the entire ambient first normal closure,
not just individual conjugates. Since \(K\le\mathcal B\),
all its retained translates are included.

The cyclic carrier \(\mathcal D\) of B.146 lies in \(\mathcal B\);
therefore
\[
 \mathcal D\cap M_i=\mathcal D\cap R_i.
\]
This does not identify that intersection with B.146's internal
cyclic relation subgroup.

\emph{B.151.4. The exact opposite-family projection.}
The left color deletion kills \(CaC,CBC,CdC,CeC\).
The actual identities
\([C,h]=1\), \(CkC=hkh^{-1}\), and \([h^2,k]=1\) give
\[
 \rho(CfC)=h(CkC)^{-1}=k^{-1}h.
\]
Since \(\theta p\) is the identity on the original region,
\eqref{eq:B151-words} gives, for every \(u\in K\),
\[
 \theta(\mathfrak d(u))
       =aC(a^{-1}\eta(u)a)Ca^{-1}.
\]
The actual \(\eta\) reverses \(f\)-exponent, so
\begin{equation}\label{eq:B151-color}
 \rho(\theta(\mathfrak d(u)))
       =a(k^{-1}h)^{-\chi_f(u)}a^{-1}.
\end{equation}
In particular all four opposite leaf cores project to one,
whereas \(\rho(\theta(\gamma))=A_0=ah^{-1}ka^{-1}\).
Writing \(k'=p(k)\), the complete opposite projections are
\begin{equation}\label{eq:B151-opposite-square}
 \rho(\theta(P_i))=1\ (1\le i\le4),\qquad
 \rho(\theta(P_0))=[(A_0k')^2,k']=:b_0.
\end{equation}

In the genuine semidirect product \eqref{eq:B151-main},
rewrite each of the five whole opposite words uniquely as
kernel times base. The results have the form
\[
 \theta(P_i)=\omega_i\ (1\le i\le4),\qquad
 \theta(P_0)=\omega_0\,\overline b_0,
\]
where the bars denote actual images in \(\mathcal B_i\).
Thus the complete initial quotient is exactly
\begin{equation}\label{eq:B151-complete}
 V_*=
 \bigl(F(\Omega_i)\rtimes\mathcal B_i\bigr)/
 \langle\!\langle
 \omega_0\overline b_0,\omega_1,\omega_2,\omega_3,\omega_4
 \rangle\!\rangle.
\end{equation}
The \(\omega_i\) are the actual coset rewrites; they are not
asserted trivial. The normal closure in \eqref{eq:B151-complete}
retains every ambient conjugate of both families. No base
section after this second filling follows. The original
fixed-source comparison and the complete cyclic-carrier
restriction remain separate questions.

\smallskip\noindent
The full precursor proof is
\path{work_DE/next_marginal_source/one_sided_family_centralizer_extension30.md},
independently reviewed in full by A, CE and root. The present
proof includes its entire presentation comparison; the record
is provenance rather than a substitute for the proof.
\end{proof}
